\section{Theoretical Introduction}

This document is an internal technical report designed to support day-to-day development, debugging, and methodological traceability of the resting-state EEG pipeline. The analytical focus remains the same as in previous versions of the project: sensor-level static functional connectivity with weighted Phase Lag Index (wPLI), estimated separately for eyes-open (EO) and eyes-closed (EC) conditions, and validated with phase-randomized surrogate inference.

The purpose of this notebook is not publication writing but operational clarity. For that reason, each section is organized around implementation decisions, data flow between stages, and reproducible execution checkpoints. All descriptions are aligned with the \Julia{} codebase and with the actual artifacts generated during runs (figures, tables, logs, and serialized intermediate files), so that any stage can be rerun in isolation when parameters, quality controls, or assumptions need to be revisited.

The report preserves the original technical depth and pipeline logic, while adopting a cleaner structure for internal work: project/data context, methodological rationale, stage-by-stage pipeline specification, algorithmic notes, and results tracking with pending actions. No conclusion section is included by design, because this notebook is intended to remain open and iterative during development.

\subsection{Multiple Sclerosis: clinical and methodological motivation}

Multiple Sclerosis (MS) is a chronic inflammatory and demyelinating disease of the central nervous system in which immune-mediated processes alter myelin integrity, axonal conduction, and glial support, with downstream effects on large-scale communication across cortical and subcortical pathways \citep{Mestre2021}. From a systems perspective, the impact is not only focal: lesions and diffuse microstructural changes can perturb synchronization, conduction delays, and compensatory recruitment patterns.

Within this context, scalp EEG is used as a millisecond-resolution window into oscillatory network behavior \citep{Cohen2014,Sanei2022}. The objective in this report is not lesion localization, but characterization of functional consequences through condition-dependent changes in spectral organization and inter-channel coupling (EO vs EC), under a reproducible preprocessing and inference framework.

\subsection{Brain Connectivity}

Resting-state EEG is commonly used to characterize spontaneous oscillatory dynamics in the absence of explicit task demands, both in healthy populations and in disease settings \citep{Cohen2014,Yamada2018}. In this context, connectivity analysis aims to quantify how signals covary or synchronize across channels and frequency ranges.

\subsubsection{Conceptual levels of connectivity}

Brain connectivity can be organized into three complementary levels. Structural connectivity refers to anatomical wiring (white-matter tracts and synaptic pathways) and constrains which regions can interact. Functional connectivity quantifies statistical dependencies between signals, without explicitly imposing directionality, and is typically estimated from covariance, spectral, or phase-based relationships. Effective connectivity goes one step further by modeling directed or causal influence from one region/signal to another under specific model assumptions.

This report emphasizes \emph{functional connectivity} at sensor level, because it is directly estimable from resting-state EEG under robust preprocessing and inferential control \citep{Cohen2014,Vinck2011WPLI,Cao2022EEGConnectivity}.

\subsubsection{Time/frequency-domain formulation}

Let $x(t)$ and $y(t)$ be two EEG signals. Their cross-correlation is
\[
R_{xy}(\tau)=\mathbb{E}[x(t)y(t+\tau)],
\]
with sample estimator
\[
\hat{R}_{xy}(\tau)=\frac{1}{N}\sum_{k=1}^{N}x_k\,y_{k+\tau}.
\]
In frequency domain, with $X(f)=\mathcal{F}\{x(t)\}$ and $Y(f)=\mathcal{F}\{y(t)\}$, the cross-spectrum is
\[
S_{xy}(f)=X(f)\,Y^*(f),
\]
where $(\cdot)^*$ denotes complex conjugation \citep{Cohen2014,Ruth2020,Sanei2022}.

\subsubsection{Non-directed connectivity metrics}

Common non-directed measures include:
\begin{align*}
\mathrm{coh}_{xy}(f) &= \frac{|S_{xy}(f)|}{\sqrt{S_{xx}(f)\,S_{yy}(f)}},\\
\mathrm{PLV} &= \left|\frac{1}{N}\sum_{k=1}^{N} e^{\mathrm{i}\Delta\phi_k}\right|,\\
\mathrm{AEC} &= \mathrm{corr}(A_x(t),A_y(t)).
\end{align*}
Coherence is amplitude+phase based and sensitive to zero-lag effects; PLV focuses on phase consistency; AEC focuses on envelope cofluctuation \citep{Cohen2014,Vinck2011WPLI,Cao2022EEGConnectivity}.

For this project, the main estimator is wPLI:
\[
\mathrm{wPLI}_{ij}
=
\frac{\left|\mathbb{E}\left[\operatorname{Im}(z_i\overline{z_j})\right]\right|}
{\mathbb{E}\left[\left|\operatorname{Im}(z_i\overline{z_j})\right|\right]+\epsilon},
\]
which downweights near-zero-lag components and is therefore better aligned with volume-conduction-robust sensor-space analysis \citep{Vinck2011WPLI}.

\subsubsection{Effective connectivity (brief theoretical frame)}

Although not the primary focus here, directed approaches are relevant as theoretical context. Granger causality evaluates whether past values of one signal improve prediction of another under autoregressive modeling; DTF and PDC quantify directed influence in frequency domain within MVAR formulations; and DCM frames directed coupling through an explicit generative model. In general, these effective-connectivity methods are more assumption-heavy than functional-connectivity metrics and require stronger model specification and validation \citep{Cohen2014,Sanei2022}.

\subsubsection{Information-theoretic measures}

Information-based alternatives include mutual information and transfer entropy:
\[
I(X;Y)=\sum_{x,y}p(x,y)\log\frac{p(x,y)}{p(x)p(y)},
\]
\[
T_{X\to Y}=\sum p(y_{t+1},y_t,x_t)\log\frac{p(y_{t+1}\mid y_t,x_t)}{p(y_{t+1}\mid y_t)}.
\]
They are flexible and potentially non-linear, but more data-hungry and computationally demanding \citep{Blinowska-Cieslak2022,Cohen2014}.

\subsubsection{Static vs dynamic connectivity}

Static connectivity summarizes one matrix over long segments:
\[
C_{\text{static}} \in \mathbb{R}^{N\times N}.
\]
Dynamic connectivity estimates time-varying matrices:
\[
C(t)=f(\text{signal}(t)),\qquad \{C_1,C_2,\dots,C_T\},
\]
typically with sliding windows. In this report, the main analysis is static to match EO/EC block design and stabilize inferential comparisons.

\subsubsection{Why this report uses CSD + wPLI + surrogates + FDR}

The selected chain is theoretically motivated and operationally coherent: optional CSD reduces zero-lag spread when active; wPLI emphasizes non-zero-lag phase consistency; phase-randomized surrogates provide empirical null distributions; and FDR controls edge-wise multiplicity. Therefore, output matrices are interpreted as statistically controlled synchronization patterns (condition/group dependent), not direct anatomical pathways or causal directions \citep{Cohen2014,Blinowska-Cieslak2022,Vinck2011WPLI,Cao2022EEGConnectivity}.
